\documentclass[../main.tex]{subfiles}
\begin{document}

%========================================================================================
% FILENAME: 05_costliness.tex (v1.0.0)
%
% §5 Costliness.  §5.1 the costliness problem (conservation identity, zeta-independence
% setup); §5.2 instantaneous costliness (proof of SP4) + attestation-append-only (the
% abstract half of SP2); §5.3 cost decomposition; §5.4 addresses and splitting.
%
% Proof budget: TWO body proofs — conservation-identity and the §5.2 proof of
% thm:interface:instantaneous-costliness.  All other proofs are \proofsketch margins.
% First-use index qualifier carried once, in the §5.2 proof.  Silences applied; cycle
% only; Douceur cite rephrased (no "Sybil"); "puppet" -> "secondary address".
%
% Change history lives in version control (see the versioning policy in main.tex).
%========================================================================================

\section{Costliness}
\label{sec:attestation:costliness}

Every claim built on $\attest{}$ requires the burn to have been costly. This chapter proves
when and how the cost arises: instantaneously and per-unit from the first cycle
(\zcref{sec:costliness:instantaneous}), but in \emph{net economic} terms only once a
market exists (\zcref{sec:costliness:problem}). Following the resource-constraint
formulation of identity multiplication\footcite{douceur2002sybil}, the burn ties standing
to a scarce redeemable asset.

% ═══════════════════════════════════════════════════════════════════════
\subsection{The Costliness Problem}
\label{sec:costliness:problem}
% ═══════════════════════════════════════════════════════════════════════

\begin{theorem}[Conservation Identity]
    \label{thm:costliness:conservation-identity}
    Let $\reservepool{} > 0$, $\totsupply{} > 0$, and
    $\ratefloor{} = \reservepool{}/\totsupply{}$. A holder with holdings $\holdings{}$ burns
    $x \leq \holdings{}$, removing those units from supply while leaving $\reservepool{}$
    unchanged. The decrease in the holder's redeemable position value is
    \begin{equation}
        \netcost
        = \frac{\holdings{}\reservepool{}}{\totsupply{}}
        - \frac{(\holdings{} - x)\reservepool{}}{\totsupply{} - x}
        = \frac{\reservepool{}\,x\,(\totsupply{} - \holdings{})}
               {\totsupply{}(\totsupply{} - x)}.
        \label{eq:costliness:conservation-identity}
    \end{equation}
\end{theorem}
\begin{proof}
    Pre-burn value at floor $\ratefloor{}$ is $\holdings{}\reservepool{}/\totsupply{}$;
    post-burn the supply is $\totsupply{} - x$, the pool is unchanged, and the holdings are
    $\holdings{} - x$, giving value $(\holdings{}-x)\reservepool{}/(\totsupply{}-x)$.
    Subtracting and collecting yields the right-hand side.
\end{proof}

\begin{corollary}[Supply-Fraction Form]
    \label{cor:costliness:supply-fraction-form}
    With $\supplyfrac{} = \holdings{}/\totsupply{}$,
    \begin{equation}
        \KEY{\;
        \netcost
        = \frac{\reservepool{}\,x\,(1 - \supplyfrac{})}{\totsupply{} - x}
        \;}
        \label{eq:costliness:supply-fraction-form}
    \end{equation}
    At $\supplyfrac{} = 0$: $\netcost \approx x\ratefloor{}$ for small $x$. At
    $\supplyfrac{} = 1$: $\netcost = 0$.
\end{corollary}

\begin{theorem}[Structural Impossibility of Internal Costliness]
    \label{thm:costliness:structural-impossibility}
    \proofsketch{$a+b = \holdings{}$ regardless of how $a,b$ are defined; $\supplyfrac{}$
    depends only on $\holdings{}/\totsupply{}$.}%
    No internal partition of a holder's holdings into sub-classes produces positive net
    economic cost at $\supplyfrac{} = 1$: the factor $(1 - \supplyfrac{})$ in
    \zcref{eq:costliness:supply-fraction-form} is independent of any such partition.
\end{theorem}

\begin{remark}[Two Regimes]
    \label{rem:costliness:two-regimes}
    Genuine net cost requires $(1 - \supplyfrac{}) > 0$ --- other holders with independent
    redemption incentives. This is a market condition, not a protocol condition; at genesis
    no competitive market exists.
\end{remark}

\begin{insight}
    \label{intuit:costliness:design-problem}
    The receipt economy is complete: instantaneous costliness, irrevocability, floor
    non-decrease, and no-recycling-advantage are protocol guarantees from the first cycle.
    Net economic costliness is a market property. The containment mechanism
    (\zcref{sec:attestation:containment}) bounds commitment capacity during the pre-market interval,
    creating the conditions under which external deposits develop the market whose feedback
    (\zcref{sec:attestation:discipline}) supplies the net-cost discipline the protocol cannot.
\end{insight}

% ═══════════════════════════════════════════════════════════════════════
\subsection{Instantaneous Costliness}
\label{sec:costliness:instantaneous}
% ═══════════════════════════════════════════════════════════════════════

\begin{proof}[Proof of \zcref{thm:interface:instantaneous-costliness}]
    \interfaceref{thm:interface:instantaneous-costliness}%
    Each increment to $\attest{a}$ is backed by \liveclass{} units redeemable at the
    settlement floor. The index $a$ is an opaque address
    (\zcref{def:interface:address-space}); the binding of addresses to any identity is the
    importing system's import-time interpretation. Burning $x$ \liveclass{} units
    (\zcref{def:operations:burn}) records on $a$ an attestation of $x\ratefloor{c}$, where
    $\ratefloor{c}$ is the floor committed by the settlement event the burn clears against;
    that floor is well-defined, and the increment order-free, under either reduction of
    \zcref{prop:invariants:burn-order-residual}.

    Credited value and destroyed value are distinct quantities. The destroyed \liveclass{}
    units were redeemable at the burn-instant floor $\ratefloor{b}$, so their value at the
    instant of destruction is $x\ratefloor{b}$. The floor is globally non-decreasing
    (\zcref{thm:invariants:floor-non-decrease}) and the settlement event precedes the burn,
    hence $\ratefloor{c} \leq \ratefloor{b}$ and
    \[
        \Delta\attest{a} \;=\; x\ratefloor{c} \;\leq\; x\ratefloor{b} \;=\; V_{\text{destroyed}}.
    \]
    Summing over the increments reaching $\attest{a} = M$ gives destroyed receipt of
    burn-instant redeemable value at least $M$ reserve units, with equality exactly when
    every burn settles at its own burn-instant floor. Strict inequality arises when an
    earlier uncleared burn has already raised the true floor above the last committed one:
    a burn of $100$ units credited at a settled floor of $1$ destroys receipt redeemable
    for $110$ when the burn-instant floor has reached $1.1$. The shortfall accrues to the
    remaining holders as floor, never to the burner as attestation, which is the
    conservative direction of \zcref{prop:interface:conservative-valuation}.

    No operation maps an attestation increment back to either class or to reserve units
    (\zcref{lem:costliness:attestation-append-only}).
\end{proof}

\begin{lemma}[Attestation Append-Only]
    \label{lem:costliness:attestation-append-only}
    \interfaceref{postc:interface:monotonicity}%
    \proofsketch{only \zcref{def:operations:burn} modifies $\attest{a}$, and only by adding
    $x\ratefloor{t} > 0$.}%
    $\attest{a}$ is monotone non-decreasing for every address $a$: no operation of
    \zcref{sec:attestation:operations} decrements it, and no operation maps the attestation back to
    either class or to reserve units. This is the abstract-layer half of
    \zcref{postc:interface:monotonicity}; the granularity corollary is established by a reader
    against its own sampling depth, not here.
\end{lemma}

% ═══════════════════════════════════════════════════════════════════════
\subsection{Cost Decomposition}
\label{sec:costliness:decomposition}
% ═══════════════════════════════════════════════════════════════════════

\begin{theorem}[$\livesplit$-Independence of Net Economic Cost]
    \label{thm:costliness:zeta-independence}
    \proofsketch{at genesis $\holdings{}+g = \livesplit\genesisreserve +
    (1-\livesplit)\genesisreserve = \genesisreserve = \totsupply{}$; $\livesplit$
    cancels.}%
    For a holder with $\holdings{}$ \liveclass{} and $g$ \tlockclass{} units, write
    $\supplyfrac{} = (\holdings{}+g)/\totsupply{}$. The conservation identity gives
    $\netcost = \reservepool{}\,x\,(1-\supplyfrac{})/(\totsupply{}-x)$; since
    $\holdings{}+g$ is independent of $\livesplit$, net economic cost is
    $\livesplit$-independent. The per-unit cost of commitment is determined entirely by
    supply dilution from external deposits.
\end{theorem}

\begin{definition}[Position Decomposition]
    \label{def:costliness:position-decomposition}
    After burning $x$ at floor $\ratefloor{}$, a holder's position separates into three
    components: the attestation obtained ($x\ratefloor{}$, non-realizable); remaining
    \liveclass{} ($(\holdings{}-x)\ratefloor{}^{+}$, immediately redeemable); and
    \tlockclass{} ($g\ratefloor{}^{+}$, model value, not realizable until maturity).
\end{definition}

\begin{definition}[Redeemable Cost Ratio]
    \label{def:costliness:redeemable-cost-ratio}
    The fraction of the burn's face value representing immediately redeemable value loss is
    \begin{equation}
        \redeemratio = \frac{\totsupply{} - \holdings{}}{\totsupply{} - x}.
        \label{eq:costliness:redeemable-cost-ratio}
    \end{equation}
    With the time-locked class, $\holdings{} = \livesplit\genesisreserve$ at genesis, so
    $\redeemratio \approx 1 - \livesplit$ for small $x$.
\end{definition}

\begin{table}[h!]
    \centering
    \caption{Redeemable cost at representative $\livesplit$. A genesis holder burns all
    $\livesplit\genesisreserve$ \liveclass{} units in infinitesimal increments.}
    \label{tab:costliness:redeemable-cost}
    \begin{tabular}{@{}c c c c c@{}}
        \toprule
        $\livesplit$ & Marginal $\redeemratio$ & $\genesisfloorfactor$
            & Net cost (certain maturity) & $\attestcap/\genesisreserve$ \\
        \midrule
        $1.00$ & $0\%$  & $\infty$ & $0\%$ & $\infty$ \\
        $0.50$ & $50\%$ & $2.00$   & $0\%$ & $0.693$ \\
        $0.10$ & $90\%$ & $1.11$   & $0\%$ & $0.105$ \\
        \bottomrule
    \end{tabular}
\end{table}

\begin{remark}[Net Cost for Supply-Dominant Holders]
    \label{rem:costliness:dominant-holder-cost}
    At $\supplyfrac{} = 1$, net cost is zero at every $\livesplit$ --- a property of closed
    reserve pools, not a deficiency. The split governs \emph{capacity}, not \emph{per-unit
    cost}.
\end{remark}

\begin{proposition}[No Recycling Advantage]
    \label{prop:costliness:no-recycling-advantage}
    \proofsketch{direct burn: $W$; infinite recycle
    $\sum_{i=1}^{\infty}\feerate(1-\feerate)^{i-1}W = W$; any mix by linearity.}%
    For fee-derived \liveclass{} units worth $W$ reserve units at the minting floor, every
    sequence of $\set{\text{burn}, \text{redeem}, \text{deposit}}$ operations produces
    operator attestation exactly $W$.
\end{proposition}

\begin{remark}[What the Time-Locked Class Does and Does Not Do]
    \label{rem:costliness:tlock-class-summary}
    \begin{enumerate}[nosep]
        \item It bounds commitment \emph{capacity} during the pre-market interval
            (\zcref{sec:attestation:containment}).
        \item Per-unit costliness holds every cycle, with or without it
            (\zcref{thm:interface:instantaneous-costliness}).
        \item Net costliness activates as external deposits dilute the genesis holder's
            supply fraction (\zcref{sec:attestation:discipline}).
    \end{enumerate}
\end{remark}

% ═══════════════════════════════════════════════════════════════════════
\subsection{Addresses and Splitting}
\label{sec:costliness:splitting}
% ═══════════════════════════════════════════════════════════════════════

The address-level analysis below is the receipt-economy form of the resource-constraint
argument: the attestation is non-transferable across addresses
(\zcref{lem:costliness:attestation-append-only}).

\begin{proposition}[Address-Splitting Cost]
    \label{prop:costliness:address-splitting-cost}
    \proofsketch{the secondary address's supply fraction $\holdings{}_P/\totsupply{}$ is
    below the funder's, so it faces higher per-unit cost.}%
    A secondary address funded by the same holder, holding $\holdings{}_P$ \liveclass{} and
    $g_P = 0$ \tlockclass{} units, burns $x \leq \holdings{}_P$:
    $\redeemratio = (\totsupply{} - \holdings{}_P)/(\totsupply{} - x)$. The attestation
    accrues to that address's $\attest{}$, not the funder's --- the attestation is
    non-transferable across addresses (\zcref{def:interface:address-space}).
\end{proposition}

\begin{proposition}[Portfolio-Level Net Cost]
    \label{prop:costliness:portfolio-net-cost}
    At $\supplyfrac{} = 1$ --- one holder controls all supply across all addresses ---
    $\netcost = 0$ regardless of the number of addresses or the distribution across them.
    The binding constraint is attestation non-transferability: $\attest{}$ on one address
    cannot be aggregated with another's.
\end{proposition}

\begin{proposition}[Splitting Invariance]
    \label{prop:costliness:splitting-invariance}
    \proofsketch{telescoping: $\totsupply{}_k - \holdings{}_k = \totsupply{} - \holdings{}$
    is invariant across sub-burns, so each sub-burn's $\netcost$ sums to the single-burn
    $\netcost$.}%
    Splitting a burn of $x$ into $n$ sub-burns preserves total redeemable cost; the total
    converges to the single-burn result as $n \to \infty$.
\end{proposition}

%========================================================================================
% END OF FILE: 05_costliness.tex (v1.0.0)
%========================================================================================

\end{document}
